\section{通过Cantor级数展开式定义实数}

记号约定:$\left(d_{n}\right)_{n=0}^{\infty}$是基列,$\left(a_{n}\right)_{n=0}^{\infty}$是$\left(d_{n}\right)_{n=0}^{\infty}$-底列.

\begin{lemma}{}{}
	设$\N$中的序列$\left(v_{n}\right)_{n=0}^{\infty}$满足$\forall n\in\N\left(0\leq v_{n}\leq a_{n}-1\right)$,$q_{0}\in\N$.对每个$m\geq 1$,令$s_{m}=q_{0}+\sum\limits_{n=1}^{m}\dfrac{v_{n}}{d_{n}}.$
	\begin{enumerate}
		\item 如果满足下列条件的实数$x$存在,那么它一定唯一,并且等于$q_{0}+\sum\limits_{n=0}^{\infty}\dfrac{v_{n}}{d_{n}}$.
		\begin{itemize}
			\item $x$在$\left(d_{n}\right)_{n=0}^{\infty}$下的Cantor级数展开式为$p_{0}+\sum\limits_{k=0}^{\infty}\dfrac{u_{k+1}}{d_{k+1}}$.
			\item $p_{0}=q_{0}$,对每个$n\in\N$有$u_{n}=v_{n}$.
		\end{itemize}
		\item 对每个$m\geq 1$有$\sum\limits_{n=0}^{\infty}\dfrac{v_{n+m+1}}{d_{n+m+1}}\leq\dfrac{1}{d_{m}}$和$\sum\limits_{n=0}^{\infty}\dfrac{v_{n+m+1}}{d_{n+m+1}}=\dfrac{1}{d_{m}}\iff\forall n\geq m+1\left(v_{n}=a_{n}-1\right)$.
		\item 如果(1)中满足条件的$x$存在,那么对每个$m\geq 1$有$s_{m}\leq x\leq s_{m}+\dfrac{1}{d_{m}}$和
		\begin{align*}
			x=s_{m}+\dfrac{1}{d_{m}}\iff\forall n\geq m+1\left(v_{n}=a_{n}-1\right)
		\end{align*}
	\end{enumerate}
\end{lemma}

\begin{theorem}{Cantor展开式的规范形式和反常形式}{proper-and-improper-expansion-of-cantor-expansion}
	设$\N$中的序列$\left(v_{n}\right)_{n=0}^{\infty}$满足$\forall n\in\N\left(0\leq v_{n}\leq a_{n}-1\right)$,$q_{0}\in\N$.令$x=q_{0}+\sum\limits_{n=0}^{\infty}\dfrac{v_{n}}{d_{n}}$,则下列情形有且仅有一种成立:
	\begin{enumerate}
		\item \textbf{情形A}:集合$\left\{n\in\N\middle|v_{n}<a_{n}-1\right\}$无限.此时,$x$在$\left(d_{n}\right)_{n=0}^{\infty}$下的Cantor级数展开式就是$q_{0}+\sum\limits_{n=0}^{\infty}\dfrac{v_{n}}{d_{n}}$.
		\item \textbf{情形B}:存在$m\in\N$使得如下条件成立:
		\begin{itemize}
			\item $\forall n\geq m\left(u_{n}=a_{n}-1\right)$.
			\item $m\geq 1\implies u_{m}<a_{m}-1$.
		\end{itemize}
		此时,存在$k\in\N$使得$x=q_{0}+\sum\limits_{n=1}^{m}\dfrac{v_{n}}{d_{n}}+\dfrac{1}{d_{m}}=\dfrac{k}{d_{m}}$,并且$\forall n\geq m+1$$\left(u_{n}=0\right)$.
	\end{enumerate}
\end{theorem}

\begin{definition}{规范展开式,反常展开式}{}
	沿用\cref{proper-and-improper-expansion-of-cantor-expansion}的设定.
	\begin{enumerate}
		\item \textbf{情形A}成立时,称$q_{0}+\sum\limits_{n=0}^{\infty}\dfrac{v_{n}}{d_{n}}$为$x$在$\left(d_{n}\right)_{n=0}^{\infty}$下的规范Cantor级数展开式.
		\item \textbf{情形B}成立时,称$q_{0}+\sum\limits_{n=1}^{m}\dfrac{v_{n}}{d_{n}}+\dfrac{1}{d_{m}}$为$x$在$\left(d_{n}\right)_{n=0}^{\infty}$下的反常Cantor级数展开式.
	\end{enumerate}
\end{definition}

\begin{definition}{终止展开式}{}
	设$x\in\R$,它关于$\left(d_{n}\right)_{n=0}^{\infty}$的Cantor级数展开式为$\lfloor x\rfloor+\sum\limits_{k=0}^{\infty}\dfrac{u_{n+1}}{d_{n+1}}$.如果
	\begin{align*}
		\exists m\in\N\forall n\geq m+1\left(u_{n}=0\right),
	\end{align*}
	则称$x$的这个Cantor级数展开式是终止的.
\end{definition}

\begin{corollary}{}{}
	设$x\in\R$,那么$x$在$\left(d_{n}\right)_{n=0}^{\infty}$下的Cantor级数展开式终止$\iff$ $x\in\bigcup\limits_{n=0}^{\infty}\dfrac{1}{d_{n}}\Z$.
\end{corollary}

\begin{proposition}{}{}
	设$x\in\Q$有既约形式$\dfrac{p}{q}$,那么下列陈述等价:
	\begin{enumerate}
		\item $x$在$\left(d_{n}\right)_{n=0}^{\infty}$下的Cantor级数展开式终止.
		\item 存在$n\in\N$使得$q\mid d_{n}$.
	\end{enumerate}
\end{proposition}

\begin{proposition}{}{}
	每个有理数$x$在$\left(d_{n}\right)_{n=0}^{\infty}$下的Cantor展开式都终止$\iff$ $\forall n\geq 1\exists n\in\N\left(q\mid d_{n}\right)$.
\end{proposition}

\begin{corollary}{}{}
	如果$\forall q\geq 1\exists n\in\N\left(q\mid d_{n}\right)$,那么对每个实数$x$都有$x\in\Q$ $\iff$ $x$在$\left(d_{n}\right)_{n=0}^{\infty}$下的规范Cantor展开式终止.
\end{corollary}

\begin{proposition}{}{evaluation-map-over-cantor-number-space}
	对每个$n\geq 1$,令$I_{n}=\left\{0,\ldots,a_{n}-1\right\}$.记$\varphi:\Z\times\prod\limits_{n=1}^{\infty}I_{n}\to\R,\left(q_{0},\left(v_{n}\right)_{n=0}^{\infty}\right)\mapsto q_{0}+\sum\limits_{n=0}^{\infty}\dfrac{v_{n+1}}{d_{n+1}}$.
	\begin{enumerate}
		\item $\varphi$是满射.
		\item 对每个$x\in\R$有$\left|\varphi^{-1}\left(\left\{x\right\}\right)\right|=
		\begin{cases}
			1,&x\notin\bigcup\limits_{n=0}^{\infty}\dfrac{1}{d_{n}}\Z,\\
			2,&x\in\bigcup\limits_{n=0}^{\infty}\dfrac{1}{d_{n}}\Z.
		\end{cases}$.
	\end{enumerate}
\end{proposition}







